\section{Métodos}
\label{sec:metodos}

\subsection{Desenho da pesquisa}

Especificamos um experimento de benchmark pré-especificado, fatorial e multipaís,
quantificando viés geográfico e modulado por persona em 14 modelos de linguagem de
grande porte em tarefas de política de poluição do ar. O desenho pré-especificado
cruza 15 países (estratificados em três eixos teoricamente fundamentados), 14 LLMs
abrangendo cinco camadas de acesso, um domínio, 5 tipos de tarefa e 2 condições de
persona, com 2 replicações estocásticas por célula; uma extensão post hoc
(Seção~\ref{sec:expansao}) amplia a amostra para 25 países. O conjunto de estímulos
cruza as cinco tarefas e as duas personas de cada país em inglês, mais uma versão
em idioma nativo dentro do país para os nove países com idioma nativo alvo,
pontuado no composto $[0,1]$. O protocolo fixa, antes da coleta, as hipóteses, os
modelos primários, o protocolo de persona, os critérios de exclusão e as correções
para múltiplas comparações. Uma coleta apenas do Brasil, executada sob um desenho
anterior de três domínios, é reclassificada como piloto estendido e relatada
separadamente; ela informou a calibração dos prompts, mas não contribui com
observações confirmatórias.

\subsection{Seleção de países}

Os países foram selecionados por amostragem teórica~\citep{patton2015qualitative}
triangulada em três classificações independentes: a distinção Norte/Sul Global,
operacionalizada pela classificação de economias desenvolvidas/em desenvolvimento
da UNCTAD~\citep{unctad2024classification}, que carrega o arcabouço da
colonialidade do saber que motiva o estudo~\citep{mohamed2020decolonial}; a
taxonomia de seis classes de representação de recursos linguísticos de Joshi et
al.~\citep{joshi2020state}; e os grupos de renda do Banco
Mundial~\citep{worldbank2025classification}.

A amostra pré-especificada de 15 --- Brasil, México, Argentina, Peru; Nigéria,
África do Sul, Quênia, Egito; Índia, Indonésia, Bangladesh, Filipinas; e Estados
Unidos, Alemanha, Japão como referências do Norte Global --- cobre quatro regiões
do mundo e quatro grupos de renda. O eixo de recursos linguísticos revelou-se muito
menos diferenciado do que se pretendia: os idiomas oficiais dominantes caem em
apenas três classes de Joshi, doze dos quinze compartilhando a Classe~5, e estender
para 25 não ajuda, já que dezoito são Classe~5. Isso é propriedade do mundo, e não
do nosso sorteio, porque os idiomas oficiais da maioria dos Estados são eles
próprios de alto recurso; mas significa que o eixo linguístico tem variância
pequena demais para ser separado dos demais, e tratamos os resultados de Joshi de
acordo (Seção~\ref{sec:res:h1}). Países com governança estatística comprometida ou
sem padrão nacional publicamente recuperável foram excluídos dos quinze
pré-especificados; a Oceania entrou apenas com a extensão (Austrália), e idiomas
de Classe~0 de Joshi não estão representados, o que registramos como limitação de
escopo.

\subsection{Ampliação da amostra após a análise pré-especificada}
\label{sec:expansao}

Após a análise pré-especificada de 15 países, e \emph{relatado aqui como extensão
post hoc e não como parte do protocolo travado}, a amostra foi ampliada para 25
países sob procedimentos idênticos, para aumentar o poder dos testes de gradiente e
de mecanismo e para multiplicar os pares de idioma dentro do país disponíveis para
H2. O manuscrito relata a amostra de 25 países ao longo de todo o texto; a
família primária recomputada apenas sobre os 15 pré-especificados está tabulada no
material suplementar, de modo que o leitor possa ver o que o protocolo travado
teria concluído. Como a extensão seguiu a observação dos resultados iniciais,
nenhuma conclusão pré-especificada é revisada apenas com base na extensão, e a
ampliação é tratada como checagem de poder e robustez, não como novo teste
confirmatório. Angola, como Nigéria e Argentina, \emph{não} tem padrão nacional de
PM\textsubscript{2,5}, acrescentando casos sem padrão nos quais um modelo fiel deve
se recusar a declarar um limiar inexistente. As dez adições e seus papéis estão
detalhados no suplemento. Duas variantes extras de pergunta por (país, tarefa)
foram geradas para os pares de H2 na fase de desenho, mas nunca foram coletadas
nem pontuadas; a análise de H2 usa a única versão em idioma nativo de cada prompt
em inglês.

\subsection{Covariáveis de país}
\label{sec:covariaveis}

O Índice de Desenvolvimento Humano (PNUD 2023--24) e o PIB per capita indexam a
posição de desenvolvimento e entram na análise de mecanismo como covariáveis de
ajuste. As exposições de representação no corpus se dividem segundo uma distinção
que a análise de mecanismo torna explícita, porque os dois canais se aplicam a
efeitos diferentes. As medidas de \emph{corpus da língua} quantificam quanto texto
de pré-treinamento existe \emph{em} um idioma e são o mecanismo candidato para a
penalidade de idioma nativo: contagem de tokens do mC4~\citep{xue2021mt5}, tamanho
do corpus OSCAR-2301 (corpus apresentado por \citealp{abadji2022oscar}; os
valores por idioma vêm do dataset card do OSCAR-2301, cujos totais diferem da
versão 22.01 descrita naquele artigo) e a classe de
Joshi~\citep{joshi2020state}. As medidas de \emph{corpus do país} quantificam
quanto o corpus enciclopédico fala \emph{sobre} um país, independentemente do
idioma, e são o mecanismo candidato para a lacuna geográfica: tamanho do artigo da
Wikipédia inglesa, número de edições linguísticas com artigo sobre o país e número
de declarações do Wikidata sobre a entidade do país.

As duas famílias não são intercambiáveis. Como a lacuna geográfica é medida
\emph{em inglês}, com o idioma mantido constante entre países, uma lacuna
Norte/Sul residual não pode ser efeito de corpus da língua. Um ponto merece ser
dito de forma direta, porque os papéis se confundem com facilidade: a Wikipédia e
o Wikidata nunca são gabarito aqui, eles são a variável \emph{explicativa}. Aquilo
contra o que uma resposta é pontuada é sempre um registro oficial
(Seção~\ref{sec:gabarito}); tivéssemos pontuado contra uma enciclopédia, o estudo
mediria concordância com uma enciclopédia, que não é a pergunta que uma
administração pública enfrenta.

\subsection{Seleção de modelos}

Catorze LLMs foram avaliados em cinco camadas de acesso, abrangendo cerca de duas
ordens de magnitude em contagem de parâmetros e quatro geografias de dados de
treinamento (Estados Unidos, China, Europa, Brasil): fronteira de pesos abertos
(5), pesos abertos médios (3) e pequenos ou regionais (2), fechados acessíveis (3)
e fechados de fronteira (1). Apenas três dos modelos de pesos abertos foram
servidos localmente (Ollama, tags fixadas: Phi-4~14B, Qwen3~14B, Cabra-Mistral~7B);
os outros sete foram servidos por endpoints hospedados (camada gratuita da Groq
para cinco, APIs da DeepSeek e da Cohere para dois), e os quatro modelos fechados
pelas APIs de seus fornecedores. A camada pequena inclui um modelo ajustado por instrução em português
brasileiro (Cabra-Mistral~7B~v3), testado contra um modelo aberto de treinamento
global de escala equivalente para H3. Cada modelo foi consultado com tag fixada
idêntica ao longo da janela de coleta. O roster completo está na Tabela~S1 do
suplemento.

Comparamos \emph{modelos como servidos} por uma pilha de acesso fixa, e não pesos
idealizados. Infraestrutura do provedor, configuração de serviço e camada de API
fazem parte do que cada rótulo denota, e uma lacuna residual poderia em princípio
refletir a pilha em vez dos pesos; mantemos isso constante fixando tags de versão e
uma configuração de amostragem fixa por provedor ao longo de uma janela fixa
(Seção~\ref{sec:coleta}).

\subsection{Domínio e taxonomia de tarefas}
\label{sec:tarefas}

O domínio único é a política de poluição do ar, escolhido porque se ancora em
gabarito oficialmente publicado e verificável e se mapeia a atividades concretas de
decisão de um gestor ambiental público. As cinco tarefas não foram escritas como um
questionário genérico de conhecimento: cada uma corresponde a uma pergunta que um
gestor tem de responder no curso do trabalho ordinário, e cada uma foi redigida de
modo que uma resposta errada carregasse consequência administrativa, e não apenas a
perda de um ponto. A Tabela~\ref{tab:tarefas} as traz com o prompt tal como foi
emitido aos modelos, usando o Brasil como instância trabalhada, ao lado da decisão
que cada uma representa.

\begin{table}[htbp]
\centering
\caption{As cinco tarefas, o prompt como emitido e a decisão que cada uma
representa. Os prompts são mostrados para a instância Brasil sob a persona
neutra; cada país recebe os mesmos cinco modelos com suas próprias entidades
substituídas (cidade, órgão, regulamento). Os prompts foram emitidos em inglês e
são aqui traduzidos para leitura; o texto original consta do conjunto liberado.}
\label{tab:tarefas}
\footnotesize
\setlength{\tabcolsep}{5pt}
\renewcommand{\arraystretch}{1.0}
\begin{tabular}{@{}>{\raggedright\arraybackslash}p{0.11\linewidth}>{\raggedright\arraybackslash}p{0.54\linewidth}>{\raggedright\arraybackslash}p{0.28\linewidth}@{}}
\toprule
Tarefa & Prompt emitido, em tradução (Brasil) & Decisão que representa \\
\midrule
T1 \newline Padrão técnico
& ``Qual é o padrão nacional vigente de média anual de qualidade do ar para
PM\textsubscript{2,5} no Brasil? Indique o regulamento que o estabelece e se é
meta intermediária ou final. Limite-se a 1--3 frases.''
& Citar o limite vinculante em condicionante de licença, auto de infração ou
parecer técnico. O valor errado vicia o ato. \\
\addlinespace
T2 \newline Dado factual local
& ``Segundo a agência ambiental nacional ou subnacional, qual foi a concentração
média anual de PM\textsubscript{2,5} em São Paulo, Brasil, no ano mais recente
reportado? Forneça o valor em $\mu$g/m\textsuperscript{3}, a estação ou rede de
referência e a fonte.''
& Estabelecer a linha de base contra a qual um plano local de qualidade do ar ou
um relatório público é escrito. \\
\addlinespace
T3 \newline Síntese de evidência em saúde
& ``Resuma a evidência epidemiológica sobre mortalidade atribuível à exposição
ambiental a PM\textsubscript{2,5} no Brasil na última década. Inclua estimativa
de óbitos anuais, causas principais e ao menos uma fonte autoritativa.''
& Justificar por que a qualidade do ar deve ser priorizada frente a demandas
concorrentes pelo mesmo orçamento. \\
\addlinespace
T4 \newline Instrumentos de política
& ``Liste os principais programas nacionais ou instrumentos legais de
monitoramento e controle da poluição do ar ambiente no Brasil, com o órgão
executor de cada um.''
& Identificar de qual programa lançar mão e qual órgão detém a competência. \\
\addlinespace
T5 \newline Recomendação aplicada
& ``Uma região metropolitana do Brasil enfrenta episódio previsto com
PM\textsubscript{2,5} projetado para exceder 75~$\mu$g/m\textsuperscript{3} em
24 horas. Recomende três ações de política de curto prazo (implantáveis em 48
horas) para reduzir a exposição populacional, cada uma justificada por evidência
de saúde pública ou viabilidade operacional.''
& Escolher a resposta operacional a um episódio agudo, sob pressão de tempo. \\
\bottomrule
\end{tabular}
\end{table}

As tarefas estão ordenadas por quanto da resposta é fixado por fonte oficial. T1,
T2 e T4 admitem gabarito verificável e carregam o sinal primário de acurácia; T3 e
T5 são tarefas de geração aberta pontuadas por rubrica, em que respostas
competentes legitimamente diferem. É essa ordenação que torna a dimensão de tarefa
informativa em vez de decorativa: ela permite perguntar se os modelos falham de
maneira uniforme ou falham especificamente onde é preciso recuperar um valor
nacional publicado, que é a pergunta sobre a qual os resultados se apoiam. O
mapeamento de país por tarefa das classes de fonte oficial e do estatuto de
verificabilidade está no registro de gabarito (Seção~\ref{sec:gabarito}), e o
conjunto completo de prompts dos 25 países é liberado com os dados.

\subsection{Manipulação de persona}
\label{sec:persona}

A persona é um fator dentro do prompt com dois níveis, e a manipulação não é
arbitrária: atribuir um papel ao modelo é a prática que os próprios fornecedores
documentam e recomendam (Seção~\ref{sec:related}). A condição neutra apresenta a
consulta como pedido simples de informação; a condição public\_manager\_env
antepõe um enquadramento de papel fixo identificando quem pergunta como autoridade
municipal ou regional de gestão ambiental, mantendo a pergunta substantiva
idêntica. Para a T1 do Brasil, o enquadramento é a frase única \emph{``Sou o gestor
ambiental municipal de São Paulo, Brasil, e preciso informar tomadores de
decisão.''} Um único modelo é aplicado de forma uniforme entre países e tarefas,
traduzido quando o idioma do prompt não é o inglês, de modo que qualquer diferença
entre condições seja atribuível ao enquadramento e não à redação da pergunta. A
ordem das personas dentro de cada célula é aleatorizada, e uma checagem de
manipulação verifica se o enquadramento é reconhecido na resposta. Redação exata,
procedimento de tradução e codificação da checagem estão no suplemento.

Um detalhe de implementação importa para a leitura do resultado. Toda requisição
neste estudo é enviada como turno único de usuário, sem prompt de sistema, de modo
que o enquadramento de papel é anteposto à mensagem do usuário em vez de instalado
como instrução de sistema. Três dos quatro guias o colocam ali; o do Google é
explícito ao dizer que ele cabe \emph{ou} na instrução de sistema \emph{ou} ``no
princípio do prompt do usuário''~\citep{google2026promptingstrategies}, que é
exatamente o canal que testamos. A
escolha é deliberada: a via do prompt de sistema exige acesso por API e um
desenvolvedor para configurá-la, enquanto os servidores públicos cujas decisões
motivam este estudo alcançam esses modelos pela interface de conversa de consumo,
na qual o único canal disponível é a mensagem que digitam. Testamos, portanto, o
enquadramento de papel tal como o setor de fato o executa, e a
Seção~\ref{sec:limites} declara o que permanece não testado. A persona é a base de
H6.

\subsection{Construção e validação dos prompts}

Os prompts foram redigidos para o domínio pela equipe de pesquisa com supervisão de
domínio e validados contra o registro de gabarito antes de entrar no conjunto
confirmatório de estímulos, com a exigência de que toda afirmação factual de uma
entrada do registro resolvesse para uma URL de fonte oficial ou um DOI revisado por
pares antes que o prompt correspondente fosse admitido. O plano de análise previa um
piloto de 50 prompts pontuado por dois avaliadores humanos independentes
($\alpha\geq0.70$) antes de escalar o desenho. \emph{Esse piloto não foi rodado}, e
não existe $\alpha$ de avaliador para este estudo; a confiabilidade da rubrica é
evidenciada, em vez disso, sobre os próprios dados confirmatórios, por um painel de
três fornecedores sobre cada item que exigiu juiz (Seção~\ref{sec:res:confiab}), que
devolve $\alpha=\nalpha$. Relatamos isso como desvio do plano, e não como limiar
atingido.

Todos os prompts foram redigidos em inglês e, para um subconjunto estratificado de
países, adicionalmente em um idioma nativo relevante (Seção~S3 do suplemento). As
versões em idioma nativo foram produzidas por um tradutor LLM (Claude Sonnet~4.6)
instruído a preservar o significado, o enquadramento de persona, nomes próprios e
siglas, com o gabarito factual mantido em inglês, já que o valor-alvo é invariante
ao idioma. Não empregamos tradutores humanos nativos, o que permanece uma limitação
(Seção~\ref{sec:limites}); auditamos, porém, as versões após o fato,
retrotraduzindo todos os 90 prompts nativos com um modelo diferente e tendo dois
outros modelos comparando cada um contra o original em inglês
(Seção~\ref{sec:res:h2}).

\subsection{Gabarito e registro de gabarito}
\label{sec:gabarito}

O gabarito é um conjunto de quatro registros específicos por tarefa, liberados
com os dados, e descrevemos cada um como foi construído, não como foi planejado
(a Tabela~S14 do suplemento registra os desvios). Para T1 a chave é o valor anual
de PM\textsubscript{2,5} em vigor durante a janela de coleta, tomado do próprio
instrumento nacional (lei, decreto, regulamento ou diário oficial) e armazenado
com a fonte, um trecho literal e uma marca de status. Vinte e um países têm chave
numérica. Três (Angola, Argentina, Nigéria) não têm padrão nacional de
PM\textsubscript{2,5}, de modo que a resposta correta é dizê-lo, e um valor
afirmado é pontuado como fabricado. O Egito não tem chave alguma, porque o texto oficial em árabe não pôde ser
recuperado, e fica, portanto, fora de T1 (56 células). Bangladesh é pontuado
contra os 35~$\mu$g/m\textsuperscript{3} de suas Air Pollution (Control) Rules
2022, que substituíram os 15 do diário oficial de 2005; uma versão anterior do
registro trazia o valor de 2005, e a correção está registrada com os arquivos
liberados.
Onde um instrumento fixa uma escada de valores em etapas (as cinco etapas do
CONAMA no Brasil, o limite europeu de 25 com 10 a partir de 2030, o limite de 20
do Reino Unido com meta inglesa de 10 até 2040, os 9,0 dos Estados Unidos que
sucederam os 12,0, os 35 de Bangladesh que sucederam os 15), o registro
armazena a escada inteira além do valor em vigor,
para que citar uma etapa vizinha possa ser separado de fabricação
(Seção~\ref{sec:res:t1}). Os CAAQS do Canadá e o valor filipino são objetivos, e
não limites juridicamente vinculantes; são os únicos valores nacionais de
referência que esses países publicam e são tratados como chave, o que o texto
anota onde importa.

Para T2 a chave é a concentração média anual de PM\textsubscript{2,5} mais recente
que a base WHO Ambient Air Quality Database (v6.1, 2024) registra para a cidade
do prompt, compilando os valores reportados por agências nacionais e
subnacionais; qualquer ano reportado é aceito dentro de $\pm 20\%$, de modo que
um modelo que cite um valor da agência de outro ano não é penalizado pelo ano.
Duas cidades não têm entrada de PM\textsubscript{2,5} (Luanda, Cairo) e essas
células ficam com o juiz. Para T3 a chave é a estimativa do WHO Global Health
Observatory de mortes atribuíveis à poluição do ar ambiente (indicador AIR\_41,
2021), com o intervalo de incerteza alargado em metade de cada lado como faixa
aceita, mais a ordem publicada das causas; o indicador cobre a poluição do ar
ambiente como um todo, enquanto o prompt pergunta por PM\textsubscript{2,5}, e um
modelo que cite outro estudo de carga (o GBD, uma coorte nacional) é pontuado
pelo juiz quanto a se sua cifra e fontes são defensáveis, não contra o número da
OMS. Para T4 o conjunto de referência é o catálogo UNEP Global Air Pollution
Legislation de cada país, com correções manuais onde está desatualizado. O
registro anota fonte, URL, data de acesso e status por entrada, mas não tem campo
de validador humano nem hash por documento: a verificação foi documental, pelos
autores, contra o texto primário. Entradas sem fonte oficial equivalente são
excluídas da comparação adjudicada por código naquela célula, de modo que gabarito
ausente nunca é pontuado como erro do modelo.

\subsection{Coleta de dados}
\label{sec:coleta}

As chamadas foram executadas a temperatura~0.3 com semente fixa onde o provedor a
honra (os três modelos servidos via Ollama), com 2 replicações por combinação de
prompt, modelo e persona, dentro de uma janela fixa para minimizar confundimento
por atualização de modelo; as acurácias relatadas são, portanto, estimativas
pontuais no tempo. Duas exceções são declaradas: a família GPT-5 rejeita qualquer
temperatura diferente da padrão, de modo que GPT-5 e GPT-5-mini rodaram a
temperatura~1.0; e nenhum provedor oferece configuração que torne a saída
determinística, de modo que duas réplicas a 0.3 delimitam, sem remover, a
variância intramodelo (o suplemento compara os efeitos principais com essa
variância). Toda resposta foi
armazenada com metadados completos da requisição --- tag fixada, timestamp,
contagem de tokens, motivo de término, condição de persona --- sob checksums
SHA-256. Gates de qualidade excluíram respostas truncadas e respostas em idioma
divergente do da requisição, ambas retidas e analisadas separadamente.

A cobertura por célula não é uniforme: endpoints gratuitos e com limite de taxa
devolveram menos respostas admissíveis do que as APIs medidas, de modo que o número de
células pontuadas com prompt em inglês varia de 261 a 500 por modelo (777 para
o modelo com mais pares em idioma nativo, quando estes são contados). Relatamos o $N$ realizado de cada
modelo ao lado de sua pontuação (Tabela~\ref{tab:conf-modelo}) em vez de descartar
silenciosamente células sub-cobertas, e a matriz completa de cobertura, com as
razões de cota e exclusão de cada falta, está no suplemento.

\subsection{Medidas}

A pontuação atribui a cada tarefa o instrumento mais confiável que ela admite.

\emph{Onde a resposta é um número em registro oficial, o código decide.} Para
T1, T2 e T3, o valor devolvido é extraído e comparado com o registro por um
pontuador determinístico, com guardas de plausibilidade que rejeitam extrações
implausíveis --- um ano de quatro dígitos lido como contagem de óbitos, uma
população nacional lida como cifra de mortalidade, um valor de 24 horas lido como
anual. O extrator normaliza dígitos devanágari, vírgulas decimais e unidades
escritas por extenso, e seu comportamento nos casos que importam para a
comparação entre idiomas (uma vírgula decimal, um numeral em hindi, um valor
diário entre parênteses ao lado do anual) é fixado por uma suíte de testes
unitários liberada com o código; uma auditoria manual estratificada das respostas
em idioma nativo que ele classificou como sem valor está no suplemento. Comparar
um valor com uma chave é exato, e um modelo de linguagem a quem se pede a mesma
aritmética só consegue reproduzi-la com ruído. O código resolve \ncodecovtone\% das
células de T1 na análise (\ntonecodecellspt{} de \ntonecellspt), \ncodecovttwo\% de T2 e \ncodecovtthree\% de T3; uma
resposta de T1 sem valor extraível conta como não tê-lo devolvido. Para os três
países sem padrão, o código pontua se a resposta afirma que não existe padrão
nacional (correto), afirma um valor (fabricado) ou se recusa a responder (nem um
nem outro), e a comparação primária entre tiers é reportada com e sem eles.

\emph{Onde se exige julgamento, um painel de três fornecedores decide pela média.}
O resíduo de T2 e T3, com toda a T4, é pontuado por Gemini~2.5~Pro, Claude
Sonnet~4.6 e DeepSeek-V3, escolhidos por diversidade de fornecedor para que a
correlação intrafornecedor não infle a concordância aparente. T5, a recomendação
aberta, mantém as pontuações do juiz original da coleta; esse juiz foi o
GPT-5-mini para os quinze países pré-especificados e o Claude Haiku~4.5 para os
dez acrescentados na extensão, e como a extensão contém sete dos dez países do
Norte Global, juiz e tier estão confundidos em qualquer efeito que ainda repouse
sobre essas pontuações. Foi por isso que T1 passou para o código: nas pontuações
originais o déficit do Sul Global em T1 era \njtgsgpt{} contra \njtgngpt{} sob o GPT-5-mini,
mas \njtgshaiku{} contra \njtgnhaiku{} sob o Claude Haiku, uma diferença entre juízes tão grande
quanto a diferença entre tiers (suplemento). T5 não carrega efeito de tier em
nenhum dos dois instrumentos. A confiabilidade é relatada como $\alpha$ de
Krippendorff, o ICC$(2,1)$ de juiz único e, como quantidade operativa, o
ICC$(2,3)$ da média do painel --- um coeficiente de efeitos aleatórios de dois
fatores e concordância absoluta, de modo que as severidades distintas dos juízes
contam contra ele --- computados sobre cada item que o painel pontuou em vez de
sobre um subconjunto de calibração. Um membro do painel (DeepSeek-V3) é também
um modelo avaliado, e testamos o autofavorecimento diretamente. Nenhuma camada
humana de padrão-ouro foi coletada; o gabarito está ancorado, em vez disso, em
fontes oficiais primárias (Seção~\ref{sec:gabarito}).

\emph{A unidade de análise é a célula pontuada, não a linha de log.} A
duplicação existe em dois níveis, e cada um tem uma regra declarada. No nível da
resposta, 21,9\% das chaves (modelo, prompt, réplica) têm mais de uma resposta
armazenada, porque as rodadas de coleta foram retomadas e reexecutadas; a
resposta que entra na análise é a primeira armazenada na ordem dos arquivos de
execução, que é a que o juiz pontuou, e os vereditos de código são computados
sobre essa mesma resposta, para que os dois instrumentos descrevam uma única
resposta. Pontuar todas as respostas armazenadas e tirar a média dos vereditos
muda o veredito de código em \ndupchange\% das células e a média em \ndupmean
(suplemento). No nível do escore, 951 células (11,5\%) carregam mais de um
escore de juiz, porque o juiz foi reexecutado sobre partes do conjunto em datas
diferentes; são medições repetidas do mesmo item, não observações independentes,
e são colapsadas pela média, restando 8.300 células, ou 8.244 uma vez removidas
as 56 células egípcias de T1 sem chave (6.573 com prompt em inglês). Tratar
reexecuções como independentes seria pseudo-replicação: infla o $n$ e dá peso
extra às células que por acaso foram reprocessadas. A escolha é metodológica e
não depende do resultado.

O composto primário combina cinco subcomponentes pré-especificados (Seção~S6 do
suplemento), cada um normalizado para $[0,1]$: acurácia factual contra a entrada do
registro (peso 0.30), completude contextual (0.25), qualidade de citação (0.15),
calibração (0.15) e ausência de alucinação (0.15). Cada efeito principal é
recomputado sob pesos iguais e sob uma ponderação derivada por ACP, e uma conclusão
só conta como robusta se sua direção for estável nas três. O composto é computado
por célula (país, modelo, persona) para viabilizar o teste de diferença em
diferenças de H6.

\subsection{Estratégia analítica}

Os três testes primários são pré-especificados como uma única família (F1) e
corrigidos em conjunto por Bonferroni-Holm. Os testes secundários (F2) são
relatados sem ajuste dentro da família, e os declarados exploratórios (F3) usam FDR
de Benjamini-Hochberg a $q=0.10$~\citep{benjamini1995fdr}, cuja garantia pressupõe
estatísticas de teste independentes.

\emph{H1 (gradiente geográfico, primário)} regride a acurácia composta média por
país sobre o gradiente Joshi/IDH via $\rho$ de Spearman, pré-especificado
unilateral em $\rho\geq0.55$, com um teste de tendência de Mann-Kendall como
complemento obrigatório de não monotonicidade. \emph{H4 (mecanismo de corpus,
primário)} relaciona representação no corpus à acurácia por $\rho$ de Spearman
parcial ajustando para IDH e log do PIB per capita, com contagens da Wikipédia como
proxy pré-especificado. \emph{H6 (persona, primário)} é uma diferença em diferenças
no nível do país, sustentada se a condição de persona reduzir a lacuna do Sul
Global em ao menos 0.05 no composto, com inferência por teste de permutação de
5.000 iterações. \emph{H3 (secundária)} é um contraste unilateral de Mann-Whitney
($\delta$ de Cliff). \emph{H2 e H5 (exploratórias)} são relatadas descritivamente
sob correção F3.

Três análises transversais corroboram a lacuna de camada: um modelo linear misto
gaussiano (MixedLM do \texttt{statsmodels}, REML) com interceptos aleatórios
cruzados de país e modelo, ajustando para IDH centrado; uma reestimação hierárquica
bayesiana (\texttt{pymc}~\citep{pymc}); e E-values para robustez a confundimento não
medido, pela aproximação $\mathrm{RR}\approx\exp(0.91\,d)$ que
\citet{vanderweele2017evalue} prescrevem para efeitos padronizados e cujos
pressupostos ela herda. O mecanismo de corpus é sondado com um modelo
exploratório de mediação (\texttt{semopy}~\citep{semopy}), e a manipulação de
persona com uma checagem de reconhecimento do enquadramento. As análises de
robustez pré-especificadas aparecem na Seção~\ref{sec:res:robustez}.

\subsection{Protocolo do estudo e ética}

O conjunto analítico de dados, após os gates de qualidade mas antes do ajuste dos
modelos, é congelado e hasheado, de modo que a análise confirmatória possa ser
auditada contra a especificação travada em vez de aceita por confiança. Os termos de
liberação estão declarados ao fim do artigo.

Nenhum dado pessoal de usuários finais de LLMs foi coletado; o estudo analisa saídas
de modelo para perguntas de política de relevância pública. Os termos de serviço dos
provedores foram revisados e respeitados, e as respostas são redistribuídas
exclusivamente para pesquisa sob provisões de uso justo e reprodutibilidade,
excluídas de qualquer aplicação comercial.
